%!TEX root = ../memoria.tex
% =============================================================================
%  CAPÍTULO 5 — CONCLUSIONES
%  Cierre del trabajo. 5 secciones cortas.
% =============================================================================

Este capítulo recoge la valoración global del trabajo realizado. Se revisa el
grado de cumplimiento de los objetivos planteados en el
Capítulo~\ref{CAP:INTRODUCCION}, se exponen los principales aprendizajes
obtenidos durante el desarrollo, se reconocen las limitaciones del trabajo
presentado y se proponen posibles líneas de continuación.


\section{Cumplimiento de los
objetivos\label{SEC:CONC_OBJETIVOS}}

El objetivo general del proyecto era diseñar, desarrollar y evaluar una
aplicación web accesible para facilitar la comunicación entre personas mayores y
sus familiares, aplicando los principios del diseño centrado en el usuario a lo
largo de todo el proceso. El resultado obtenido —\emph{ConectaMayor}, una
aplicación funcional y evaluada con usuarios— permite afirmar que este objetivo
se ha cumplido en lo esencial. A continuación se valora el grado de cumplimiento
de cada uno de los objetivos específicos:

\begin{enumerate}
  \item El análisis del estado del arte ha permitido situar
  \emph{ConectaMayor} en el panorama actual de las soluciones digitales
  orientadas a personas mayores y justificar las decisiones de diseño adoptadas,
  especialmente el carácter privado, accesible y acotado de la aplicación frente
  a las plataformas generalistas.

  \item La arquitectura modular se ha implementado en cinco aplicaciones Django
  independientes —usuarios, agenda, contactos, mensajes y galería— acompañadas
  del subsistema transversal de notificaciones. El sistema de tres roles
  diferencia con claridad a la persona mayor, el familiar editor y el familiar
  de solo lectura, ajustando los permisos y la interfaz a las necesidades de
  cada perfil.

  \item La interfaz simplificada del rol \emph{mayor} se ha materializado en una
  pantalla principal con cuatro accesos directos a los módulos, tipografía
  ampliada, indicadores numéricos de novedades y botones de gran tamaño. En
  paralelo, la interfaz del rol \emph{familiar editor} mantiene las funciones de
  gestión necesarias sin comprometer la simplicidad de la experiencia diseñada
  para la persona mayor.

  \item El cumplimiento de las pautas WCAG 2.1 nivel AA se ha abordado de forma
  transversal mediante contraste mínimo 4{,}5:1, áreas de interacción de al
  menos 44~$\times$~44 píxeles, indicadores de foco visibles, lenguaje sencillo
  y soporte de tres temas visuales alternativos.

  \item La privacidad se ha garantizado mediante el aislamiento de los grupos
  familiares en las consultas a la base de datos. Ningún usuario puede acceder a
  información perteneciente a un grupo familiar distinto del suyo, y la
  aplicación no integra servicios externos ni publicidad.

  \item La evaluación de usabilidad con dos personas mayores, conducida conforme
  a las recomendaciones de la Asociación Interacción Persona-Ordenador, ha
  permitido detectar doce mejoras concretas que se han incorporado a la versión
  final del sistema. Además, una limitación de diseño identificada durante la
  reflexión posterior a las pruebas —la asunción inicial de un único usuario
  mayor por grupo familiar— motivó un rediseño del modelo de datos y de la
  interfaz del editor.

  \item La presente memoria documenta de forma estructurada el proceso completo:
  desde el análisis del problema hasta la evaluación con usuarios, incluyendo
  las decisiones técnicas intermedias y su justificación. Esta documentación
  sirve tanto como registro del trabajo realizado como base para las posibles
  ampliaciones futuras descritas más adelante.
\end{enumerate}


\section{Aprendizajes obtenidos\label{SEC:CONC_APRENDIZAJES}}

Más allá del producto entregado, el desarrollo del proyecto ha generado
aprendizajes significativos. Entre ellos, destacan especialmente tres.

El primero, de carácter técnico, ha sido la consolidación práctica del
\emph{stack} Django en un proyecto de tamaño medio. Aunque durante el grado se
han tratado distintos aspectos del desarrollo web, no es lo mismo trabajar con
ejemplos guiados que enfrentarse a las decisiones reales de un proyecto iniciado
desde cero: estructura de aplicaciones, modelo de usuario personalizado, gestión
de migraciones, configuración del servidor y organización del código. La
experiencia ha mostrado el valor de tomar ciertas decisiones desde las primeras
fases —como la definición de \texttt{AUTH\_USER\_MODEL} antes de la primera
migración— y el coste que puede tener posponerlas.

El segundo aprendizaje, de carácter metodológico, ha sido la importancia de
evaluar el sistema con usuarios. Las sesiones de usabilidad han revelado
problemas que no se manifestaban durante el desarrollo, porque el propio
diseñador acaba interiorizando la lógica de su aplicación y deja de percibir
algunos puntos de fricción. Observar a una persona mayor utilizar la aplicación
durante una sesión de prueba proporciona información cualitativa que ninguna
prueba automatizada puede ofrecer. De hecho, varias de las mejoras incorporadas
en la versión final habrían pasado desapercibidas sin esta etapa.

El tercer aprendizaje, de carácter personal, ha sido constatar que la
accesibilidad no debe entenderse como un añadido cosmético al final del
desarrollo, sino como una decisión estructural que condiciona desde la
organización de la interfaz hasta la elección de palabras de cada botón. Aplicar
las pautas WCAG 2.1 ha exigido revisar componentes que en otro contexto podrían
haberse dado por válidos, y ha enseñado a analizar las interfaces digitales con
un nivel de atención mayor.


\section{Limitaciones\label{SEC:CONC_LIMITACIONES}}

El trabajo presentado tiene limitaciones que conviene reconocer explícitamente,
tanto por honestidad académica como porque marcan el punto de partida natural de
las líneas de continuación.

La primera y más evidente es que la aplicación no se ha desplegado en un entorno
de producción real. Todas las pruebas —incluidas las sesiones de evaluación con
usuarios— se han realizado sobre un servidor local accesible mediante la red
WiFi doméstica del matrimonio participante. Aunque esta configuración ha sido
suficiente para validar el comportamiento funcional del sistema, quedan fuera
del alcance del proyecto aspectos importantes como el rendimiento bajo
concurrencia, la configuración de HTTPS, la administración de un servidor público
o la gestión de copias de seguridad.

La segunda limitación afecta a la evaluación de usabilidad. El número de
participantes (n=2) es reducido para extraer conclusiones cuantitativas
generalizables, y ambos participantes comparten un perfil demográfico homogéneo:
edades cercanas, mismo nivel formativo e idéntico entorno tecnológico. Aunque
los hallazgos son útiles dentro de una evaluación cualitativa orientada a la
detección de problemas de uso, una validación más amplia requeriría diversificar
los perfiles y aumentar el número de sesiones.

La tercera limitación es técnica. SQLite, elegida como base de datos por su
simplicidad y por ser la opción por defecto de Django, resulta adecuada para el
desarrollo y para un escenario de demostración, pero no sería la opción más
conveniente para un despliegue en producción con múltiples familias accediendo
simultáneamente. Asimismo, el conjunto de pruebas automatizadas del proyecto se
limita a las funcionalidades críticas y no constituye una \emph{suite} completa,
lo que reduce la confianza ante futuras modificaciones del código sin
verificación manual.


\section{Trabajo futuro\label{SEC:CONC_FUTURO}}

A partir de las limitaciones anteriores y de las observaciones recogidas durante
la evaluación, se identifican varias líneas naturales de continuación para el
proyecto.

\paragraph{Despliegue en producción.}
La línea de trabajo futuro más inmediata sería llevar \emph{ConectaMayor} a un
servidor accesible públicamente. Esto implicaría migrar la base de datos a
PostgreSQL, configurar un servidor de aplicaciones como Gunicorn detrás de un
\emph{proxy} inverso, habilitar HTTPS mediante certificados de Let's Encrypt y
adquirir un dominio. Un despliegue real permitiría ampliar la evaluación con
familias diversas, no necesariamente cercanas al autor.

\paragraph{Notificaciones externas.}
El subsistema de notificaciones implementado actualmente se limita a indicadores
visibles dentro de la propia aplicación. Una ampliación útil sería incorporar
notificaciones por correo electrónico —para los familiares— y notificaciones
\emph{push} a través del navegador, de modo que los usuarios puedan recibir
avisos de recordatorios próximos o mensajes nuevos sin necesidad de tener la
aplicación abierta.

\paragraph{Integración con calendarios externos.}
Permitir la sincronización del módulo de agenda con servicios como Google
Calendar facilitaría que un familiar registrase una cita en su agenda personal y
esta apareciese automáticamente como recordatorio en la aplicación del mayor,
reduciendo así la duplicación de información entre sistemas.

\paragraph{Evaluación ampliada.}
Una nueva ronda de pruebas AIPO con un número mayor de participantes y con
perfiles más diversos —en edad, formación, autonomía tecnológica y contexto
familiar— permitiría comprobar hasta qué punto los hallazgos actuales se
mantienen en otros casos de uso y, posiblemente, identificar nuevos puntos de
mejora que no han emergido con el matrimonio participante.

\paragraph{Mejoras adicionales de accesibilidad.}
La accesibilidad puede llevarse más allá del nivel AA mediante la incorporación
de lectura por síntesis de voz para los textos de la interfaz, dictado por voz
para los mensajes y recordatorios, y adaptaciones específicas para usuarios con
dificultades motrices. La API \texttt{SpeechSynthesis} del navegador, ya
estandarizada, permitiría abordar la primera de estas mejoras con un esfuerzo de
implementación reducido.


\section{Reflexión final\label{SEC:CONC_REFLEXION}}

Las dos sesiones de evaluación con personas mayores han sido una de las
experiencias más valiosas del proyecto. Observar cómo una persona interactúa con
una interfaz diseñada por uno mismo —y detectar, en tiempo real, los puntos en
los que el diseño no resulta tan claro como se esperaba— proporciona un
aprendizaje difícil de obtener únicamente mediante documentación técnica o
pruebas automatizadas. También ha confirmado que el trabajo tiene sentido más
allá del marco académico: las dos personas que participaron en las pruebas
valoraron positivamente la idea y manifestaron interés en seguir utilizándola.

La brecha digital que afecta a las personas mayores no se resolverá con una sola
aplicación, pero sí puede reducirse mediante una forma distinta de diseñar
tecnología: una forma que sitúe a los usuarios reales —con sus limitaciones, su
experiencia y sus expectativas— en el centro del proceso. Si este Trabajo Fin de
Grado contribuye, aunque sea modestamente, a reforzar esa perspectiva, su valor
no estará únicamente en el código desarrollado, sino también en la manera de
entender el diseño tecnológico que ha ayudado a formar.
